% !TeX root = ../main.tex
% Included by main.tex through ch6/ch6_main.tex.
% Recovered revision v2: reconstructed from the original chapter after external directory replacement.

\providecommand{\meCDI}{\textsc{corpus-CDI}}
\selectlanguage{english}
\crefname{chapter}{chapter}{chapters}
\Crefname{chapter}{Chapter}{Chapters}

\chapter{General Discussion}
\label{chap:general-discussion}

\objectif{Synthesize what the empirical studies establish about calibrated child--model comparisons, frequency-sensitive learning, episodic retrieval, and caregiver interaction, while distinguishing behavioral resemblance, learning effects, and claims about developmental mechanisms.}

\section{Revisiting the Central Question}
\label{sec:gd_overview}

This thesis asks how limited and uneven linguistic experience can support expressive and grammatical behavior, and what language models reveal about the resources involved. It approaches this question through comparisons of expressive vocabulary, an intervention on access to prior instances, a benchmark of interaction, and a comparison of caregiver-derived rewards. These studies require a corresponding methodological question: under which learning conditions and behavioral measures can model results inform theories of language development? Computational models are valuable because they make learning assumptions explicit and allow experience, architecture, and learning signals to be manipulated under controlled conditions \citep{dupoux2018cognitive,warstadt2022what}. Their scientific interpretation is nevertheless conditional. A model can reproduce a behavioral pattern for reasons that differ from those that produce the pattern in children, and it can fail because the modeled input, task, response channel, or learning objective omits resources available to the child. Model performance becomes developmental evidence only when the comparison and the inference drawn from it are specified together.

\Cref{chap:litreview} developed this argument as a framework for calibrated child--model comparisons. It distinguished the alignment of learning experience from the alignment of observed outcomes, and it emphasized that neither endpoint similarity nor a nominally child-scale training corpus is sufficient on its own. The empirical chapters then used models in several different roles. \Cref{chp2_chap} treated language models as learners whose expressive outputs could be compared with child vocabulary production under controlled exposure and generation budgets. \Cref{chp3_chap} used retrieval augmentation as an intervention on a model's access to prior instances, asking whether that intervention selectively reduced lexical-frequency effects on syntactic contrasts. \Cref{chp4_chap} first evaluated large language models as simulators of child--caregiver dialogue and then used smaller models as controlled learners to test the effects of caregiver-derived rewards.

These studies provide two complementary forms of evidence. The \meCDI{} and dialogue studies are primarily \emph{diagnostic}: they ask where model behavior corresponds to, or diverges from, a human reference under an explicit measurement protocol. The retrieval and feedback studies are primarily \emph{interventional}: they vary a computational resource or learning signal and measure the resulting change. This distinction matters because a diagnostic resemblance does not identify the process that produced it, whereas an intervention can establish an effect inside a model without showing that the same intervention operates in children. The four studies should therefore not be read as stages of one system that first learns vocabulary, then gains episodic memory, and finally enters social interaction. They are linked by a common inferential discipline and by a set of recurring empirical questions, not by shared parameters or a common developmental timeline.

Three conclusions emerge. First, developmental comparison requires calibration of the learner's experience, the learner and intervention, the observed outcome, and the linking assumptions that connect an outcome to a theoretical construct. Second, frequency is a recurrent boundary condition: the predictive learners studied here make low-frequency lexical evidence less reliably available, while direct access to stored instances provides only partial and task-specific compensation. Third, social interaction cannot be reduced either to plausible isolated utterances or to a generic feedback signal. Sustained dialogue and different forms of caregiver response place distinct demands on a model, and their effects depend on the outcome used to evaluate them.

The resulting picture is neither a general success nor a general failure of language models as accounts of development. Predictive learning captures substantial regularity; retrieval and feedback can selectively change what a model does; and prompting can elicit some age-associated patterns. The studies also expose specific mismatches: vocabulary-growth profiles differ across learners, retrieval leaves a residual frequency gap, and improvements in generated utterances do not automatically extend to controlled grammatical preferences or sustained dialogue. No single intervention was evaluated across all these outcomes. The contribution of the thesis is thus to identify what each result licenses, which mismatches remain, and how those mismatches can be converted into more precise hypotheses for future modeling.

\section{Methodological and Computational Contributions}
\label{sec:gd_contributions}

The thesis contributes experimental tools as well as conclusions about model behavior. The first methodological contribution is a production-oriented vocabulary protocol. \meCDI{} makes three usually separate choices jointly inspectable: how much language a model encounters, how much it is allowed to produce, and which frequency distribution is represented in the evaluation vocabulary. Combining the resulting score with lexical diversity and non-word rate makes it possible to distinguish broader vocabulary coverage from a change in how widely or accurately a model samples word forms. This supplies a concrete implementation of outcome calibration from \cref{chap:litreview}, while exposing the remaining difference between a count of generated forms and a parent's report of meaningful word use.

The retrieval study contributes a controlled application and analysis of an existing AI method. It does not introduce $k$NN language modeling itself. Its contribution is to test whether interpolation with a datastore-derived distribution changes syntactic preferences selectively across lexical-frequency bands, then to vary the stored representations, structural content, granularity, context, and number of neighbors. This design moves beyond an aggregate retrieval score: it identifies conditions in which instance access helps, conditions in which it can hurt, and the information associated with the benefit. The frequency-stratified evaluation is consequently useful both for a computational hypothesis about memory and for diagnosing retrieval-augmented language models.

The interaction benchmark contributes a complementary evaluation design: it tests both interlocutor roles, measures word-, utterance-, and dialogue-level properties, and contrasts responses to human turns with sustained exchanges between model instances. The feedback study then contributes a comparative use of reinforcement learning. It holds the learner family and fine-tuning framework constant while replacing the caregiver-derived reward, and evaluates the resulting changes through two generation measures and controlled minimal pairs. Its novelty lies in the comparison of candidate signals and their differentiated effects, rather than in a new PPO algorithm. Together these contributions show how established NLP methods can be organized into informative cognitive-science experiments without requiring the four studies to share one artificial learner.

\section{From Model Performance to Developmental Evidence}
\label{sec:gd_calibration}

The studies collectively extend the two-axis framework of input and outcome alignment into a broader inferential chain:
\[
\begin{aligned}
\text{experience} &\longrightarrow \text{learner/intervention}
\longrightarrow \text{observed outcome} \\
&\longrightarrow \text{linking model}
\longrightarrow \text{licensed claim}.
\end{aligned}
\]
Each link contributes assumptions. Developmental interpretation is strongest when those assumptions are reported and varied independently, and weakest when differences at several links are compressed into a single label such as ``child-scale,'' ``human-like,'' or ``language competence.''

\subsection{Calibrating the Learner's Experience}
\label{sec:gd_calibration_input}

The amount of input is an important part of a developmental comparison, but it is not a sufficient description of experience. Children encounter speech in recurring activities, with prosody, visible events, interlocutors, and consequences for their own participation. A textual corpus can match the estimated number of words heard by a child while differing in register, modality, temporal organization, and interactional contingency. Conversely, a small naturalistic corpus may preserve child-directed register but sample only a narrow portion of the child's environment.

This distinction appears differently across the empirical chapters. In \cref{chp2_chap}, child-directed transcripts and audiobook text provide contrasting registers under exposure-scaled training conditions. The design makes nominal corpus size explicit, but successive scales do not automatically represent one continuously developing learner, and repeated optimization over the same material contributes to effective exposure. In \cref{chp3_chap}, the smaller model was trained on approximately 50 million words drawn from CHILDES, conversational transcripts, stories, and subtitles. This is a mixed BabyLM-derived corpus, not 50 million words of child-directed speech alone. The bounded-data and large-scale settings provide evidence that a retrieval effect can be observed under two different regimes; because model size, corpus composition, and training history differ simultaneously, and datastore provenance is incompletely reported, they are replications across settings rather than a controlled scale comparison. In \cref{chp4_chap}, prompted dialogue models inherit broad pretraining, whereas the feedback-learning models are pretrained on a bounded caregiver corpus and receive additional information through learned rewards.

A complete exposure account must therefore distinguish unique linguistic material from total token presentations and from externally supplied information. The ten training epochs reported for the smaller retrieval model, for example, are additional presentations of the bounded corpus, not ten independent samples of a child's experience. Retrieval datastores, prompt demonstrations, pretrained tokenizers, reward predictors, and automatic evaluators do not all function as direct input to the target learner, but they affect what the complete experimental system can do. Documenting this \emph{total information budget} requires separating resources that affect learning or inference from evaluators used only to measure the result; an evaluator does not transfer its pretraining to a learner merely by scoring it. It makes clear whether a result concerns learning from child-comparable experience, the use of a resource learned elsewhere, or an engineering system that combines both.

\subsection{Specifying the Learner and the Intervention}
\label{sec:gd_calibration_learner}

Architecture is most informative when it is treated as an experimental variable rather than assumed to be cognitively plausible by resemblance. The character-level LSTM and Transformer in \cref{chp2_chap} test whether the observed expressive-production pattern depends on one of two sequence-learning architectures. The comparison reduces dependence on a single model family, but shared failure does not isolate a unique cause: the systems also share a textual response channel, a predictive objective, and a generation-based task model. Their contrast is consequently evidence about robustness across the tested architectures, not proof of a general limitation of statistical learning.

The intervention is more direct in \cref{chp3_chap}. The parametric model and its retrieval-augmented counterpart are evaluated on the same syntactic contrasts, allowing the effect of access to stored instances to be estimated within a common outcome. The full-retrieval keys are contextual representations from the GPT-2 model itself. The interpretation is tied to the datastore, representation, interpolation rule, and retrieval configuration. The result concerns what this form of non-parametric access makes sufficient for the tested decisions; it does not establish that $k$NN retrieval is a process model of hippocampal memory.

The two studies in \cref{chp4_chap} intervene at still different stages. Prompting changes the context in which a pretrained simulator expresses its existing capacities. Reinforcement learning changes the model's production policy in response to a frozen reward. Neither manipulation should be described simply as ``more interaction.'' One supplies examples at inference time, while the other repeatedly reinforces utterances according to a learned scalar proxy. Naming the stage and content of each intervention is necessary before their effects can be compared.

\subsection{Aligning Outcomes and Task Models}
\label{sec:gd_calibration_outcome}

Children and models rarely expose the same dependent variable directly. Child knowledge is inferred from reports, looking, choices, elicited responses, or spontaneous production. A model exposes probabilities, representations, or generated sequences. A threshold, prompt, decoder, distance function, classifier, or aggregation rule connects that output to the construct under study. This additional component is a \emph{task model}, and its assumptions form part of the empirical claim.

The \meCDI{} production threshold in \cref{chp2_chap}, for example, links parent-reported vocabulary and the frequency of word forms in sampled model output. It improves comparability of observation opportunity, but it does not make the two measures identical. The syntactic minimal pairs in \cref{chp3_chap} provide a controlled test of probabilistic preference, not a direct measure of spontaneous grammatical production. The dialogue benchmark in \cref{chp4_chap} aggregates lexical, utterance, and conversational statistics, each of which captures a selected aspect of interaction rather than communicative adequacy as a whole. The feedback study then shows explicitly that two measures intended to bear on grammaticality can disagree: rewards can improve the grammaticality of generated utterances without reliably improving BLiMP or Zorro preferences.

Such disagreements should not be collapsed into an average score. They reveal that an intervention changes some expressions of a capacity but not others. Stronger claims require convergence across outcomes with different measurement biases, or a theory that predicts why a change should be restricted to one outcome. Linking functions should, where possible, be selected on data separate from the final comparison and tested for sensitivity to plausible alternatives. Otherwise, part of the apparent alignment can be introduced by the criterion designed to measure it.

\subsection{What Success and Failure License}
\label{sec:gd_calibration_claims}

An aligned success supports a bounded sufficiency statement: the specified learner can produce the specified outcome from the specified information under the specified task model. It does not establish that children use the same internal representation or optimization procedure. A controlled contrast can strengthen the conclusion by locating which manipulated resource changes performance. Thus retrieval can be said to cause a change in the evaluated model, and a reward can be said to alter the trained policy, while the developmental relevance of that change remains a separate inference.

Failure is less decisive when several mismatches remain. A text-only learner can fail because it lacks perceptual grounding, because its objective gives little weight to rare events, because its memory cannot preserve particular episodes, or because the evaluation does not measure the knowledge that changed. Such failure does not demonstrate that children require an innate constraint or that the target is unlearnable from experience. It becomes theoretically informative when alternative explanations are converted into controlled comparisons. In this sense, the residual gaps identified throughout the thesis are not merely negative results. They define the next variables that a developmental model must expose.

\section{Learning from Sparse Evidence}
\label{sec:gd_sparse_evidence}

A recurring result across the vocabulary and retrieval studies is that model behaviour becomes less robust when the relevant linguistic evidence is infrequent. This recurrence provides a genuine connection between \cref{chp2_chap} and \cref{chp3_chap}, but the connection is diagnostic rather than a direct causal chain. The two chapters examine different linguistic domains, models, and outcome measures. \Cref{chp2_chap} measures how reliably target words appear in free generation under calibrated exposure and production budgets, whereas \cref{chp3_chap} measures preferences between minimally different sentences in controlled syntactic contrasts. The shared result is therefore not that both chapters reveal a single latent deficit, nor that retrieval solves the vocabulary-production problem. Rather, they provide converging evidence that frequency is an important boundary condition on what predictive language learners make available under different forms of evaluation.

\subsection{Frequency Sensitivity in Expressive Production}

The contribution of \meCDI{} in \cref{chp2_chap} is to make model generation more comparable to observations of children's expressive production. Training exposure, generated output quantity, and the frequency distribution of the evaluation vocabulary are brought under explicit control. Within this protocol, the models differ from children not only in their final scores but also in the shapes of their trajectories. Over the evaluated exposure conditions, their generated vocabulary coverage grows more slowly and approximately linearly, and low-frequency target words are especially unlikely to reach the production criterion. Changing the sampling temperature alters the severity of this pattern, but does not yield a setting that simultaneously aligns vocabulary coverage and lexical diversity with the child data.

The discrepancy is not confined to vocabulary coverage, but neither does every model diverge on every individual measure. Lexical diversity and non-word rates vary with architecture and corpus: some CDS-trained conditions approach the child reference in type--token ratio, and the CDS-trained LSTM shows a marked decline in non-word rate. The informative mismatch concerns the joint production profile across exposure conditions. Local agreement or improvement on one measure does not establish simultaneous agreement in reliable expressive coverage, lexical diversity, and word-form accuracy. The content-word analysis adds a further comparison of how output is allocated between content and function words. Together these measures distinguish changes in the number of different forms produced from changes in how reliably evaluation words are available. The non-word contrast additionally involves different response channels: dictionary membership in orthographic output cannot isolate children's articulatory development.

The frequency-band analysis summarizes fitted curves by the exposure-equivalent point at which they cross an 80\% criterion. Some crossings lie beyond the observed training range, so they are extrapolated properties of the fitted curves rather than observed acquisition ages. Their magnitude depends on the curve and the exposure-to-age mapping. The distinction is particularly important because the CDS conditions reuse the 30-month-equivalent checkpoint for later production budgets: a larger generated sample at months 31--36 is not additional training. The supported synthesis is strong frequency sensitivity under this protocol; it does not require interpreting distant fitted crossings as a literal number of months that a learner would need.
% AUTHOR CHECK: reconcile exposure conversion and observed versus extrapolated ranges in Chapter 3 before reporting exact months-to-criterion comparisons.

This result should be interpreted at the level of the measurement actually used. A \meCDI{} score records whether a word is produced sufficiently often in a finite generated sample. It therefore measures \emph{reliable expressive availability under a specified generation protocol}; it is not a direct test of whether the model possesses a lexical representation or understands the word. Failure to cross the count threshold may reflect weak representation, low contextual probability, competition from more frequent alternatives, or the interaction between decoding and the evaluation rule. Conversely, crossing the threshold does not establish child-like semantic knowledge. The threshold is a linking function between two different response systems---parent-reported child vocabulary and sampled model text---and conclusions must remain conditional on that link.

With this qualification, the chapter nevertheless identifies an informative asymmetry. High-frequency forms occupy a disproportionately large share of model output, whereas many low-frequency content words remain difficult to elicit reliably. This pattern is compatible with the hypothesis that next-token prediction, together with the tested architectures and input distributions, allocates behavioural probability unevenly across the lexical long tail. It does not, by itself, establish cross-entropy learning as the unique cause. Corpus composition, representational choices, model capacity, repeated exposure, and decoding all contribute to the observed learner--output system. The appropriate conclusion is consequently that the tested predictive learners do not reproduce children's relative robustness to sparse lexical evidence under the \meCDI{} production measure.

\subsection{Retrieval as Selective Compensation}

\Cref{chp3_chap} asks a narrower intervention question: when parametric predictions are frequency-sensitive, can explicit access to stored instances selectively improve performance? The answer is conditionally positive. Adding $k$NN retrieval produces larger gains for low-frequency than for high-frequency items across the tested syntactic phenomena and model settings. The inferential analyses support the selectivity of this effect: retrieval is not merely a uniform increase in accuracy. At the same time, low-frequency items remain less accurate than high-frequency items after augmentation. Retrieval narrows the gap without eliminating it, and some high-frequency conditions show no benefit or modest degradation.

The size of the effect can be seen directly in the reported accuracies. For the smaller model, low-frequency subject--verb agreement improves from 0.60 to 0.76 while high-frequency agreement changes from 0.91 to 0.90. Low-frequency relative-clause accuracy rises from 0.48 to 0.65, while high-frequency accuracy changes from 0.88 to 0.89 (Table~\ref{tab:frequency_effects}). These within-model contrasts make clear that the narrower gap includes a substantial improvement for the difficult items, rather than arising only from degraded performance at the high-frequency end. The item-bootstrap analyses support the greater average gain for low-frequency items (Table~\ref{tab:boot_ci}). They do not make absolute frequency bands comparable across the two differently trained models.

The content and configuration analyses further constrain this result. Semantic similarity alone is generally insufficient to reproduce the syntactic gains, while retrieval sets containing structurally relevant instances tend to be more useful. In the low-frequency conditions, similarity-based selection without a structural match does not outperform random selection from sets containing a structural match (Figure~\ref{fig:datastore}). The relevant match is operationalized by detection of the target phenomenon, rather than identity of a complete syntactic derivation. However, this does not demonstrate separable, independently necessary structural and semantic components of human episodic memory. Contextual representations may retain several kinds of information at once, and apparently ``semantic-only'' representations may not remove structure completely. Moreover, the best retrieval granularity, context size, and number of neighbours vary across phenomena. The observed benefit is therefore produced by a particular interaction among the base model, datastore, query representation, retrieval rule, and target construction.

Accordingly, the $k$NN system should be understood as a computational intervention that isolates one useful property of an episodic system: access to particular stored instances without additional gradient updates. It does not instantiate the full set of functions attributed to human hippocampal memory, such as selective consolidation, adaptive encoding, online updating, or task-sensitive retrieval control. The results show that instance-based access is sufficient to improve some low-frequency syntactic decisions in the tested models. They do not show that the same mechanism explains children's frequency robustness or that the hippocampus is the source of that robustness.

Taken together, the two chapters motivate a plural rather than a replacement account. Parametric prediction captures substantial distributional structure but remains frequency-sensitive; non-parametric access can compensate selectively when useful instances are available and retrievable. This is a principled connection between the chapters, but it remains a hypothesis about complementary computational resources. A stronger test would apply retrieval to the same expressive-production outcome used in \cref{chp2_chap}, or compare parametric and retrieval-based learners on matched production and minimal-pair tasks. Such work would also need to count the datastore as part of the learner's total information budget, following the calibration principles developed in \cref{chap:litreview}. Retrieval does not provide information for free: its developmental interpretation depends on what is stored, when it becomes available, and how access is controlled.

\section{Competence Depends on How It Is Measured}
\label{sec:gd_measurement}

The studies in this thesis do not reveal a single quantity that can be called ``language competence.'' Instead, they show that an apparent success or failure depends on the behavioural outcome, the task used to elicit it, and the linking assumptions connecting that outcome to the construct of interest. This dependence is not merely methodological noise. It determines which developmental claims the evidence can support.

\subsection{Endpoints and Trajectories Answer Different Questions}

An endpoint comparison asks whether a model eventually reaches a specified level of performance. A trajectory comparison asks whether changes with increasing experience resemble the timing, ordering, and shape observed in children. These questions can dissociate. Two learners may reach similar endpoints through very different paths, while learners with similar early trajectories may later diverge. The vocabulary study makes this distinction especially clear: a model's eventual production of more CDI words does not compensate for a growth curve, frequency dependence, or lexical allocation that differs from the child trajectory.

Developmental labels must therefore be attached only to manipulations that warrant them. The exposure-scaled conditions in \cref{chp2_chap} approximate increasing quantities of linguistic experience, but they are not equivalent to observing one child longitudinally, particularly when separate training sets or models contribute to successive points. Likewise, the age-conditioned prompts in the dialogue study of \cref{chp4_chap} test whether a pretrained model can \emph{simulate} age-associated behaviour. They do not constitute evidence that the model learned through those developmental stages. Reinforcement learning in the same chapter is a genuine intervention on model behaviour, but its before--after contrast is still not a child developmental trajectory. Keeping these comparisons distinct prevents endpoint resemblance, prompted role performance, and learning-induced change from being treated as interchangeable evidence.

\subsection{Free Generation and Controlled Preferences Expose Different Capacities}

Free generation is valuable because model and child outputs can be analysed using related production measures. It is also jointly determined by representation, contextual preference, decoding, output length, and sampling variability. In \cref{chp2_chap}, a target word that rarely appears may be weakly represented, or it may simply lose repeatedly to more probable alternatives. In the feedback study of \cref{chp4_chap}, improvements in generated grammaticality demonstrate a change in the production policy, but do not necessarily entail a corresponding change in abstract grammatical preferences.

Minimal-pair evaluation reduces some of this ambiguity. The contrasts in \cref{chp3_chap} hold much of the lexical and propositional content constant and ask whether the model assigns greater probability to the grammatical member. This design is more sensitive to a particular distinction, but it is also narrower and less naturalistic. Success does not establish that the same knowledge will guide spontaneous production or interaction. Conversely, weak free generation does not imply failure on a controlled preference test.

The divergence between generation-based grammaticality and BLiMP/Zorro-style judgments after feedback learning in \cref{chp4_chap} is therefore theoretically informative. It suggests that an intervention can make sampled utterances more acceptable without producing measurable transfer to controlled syntactic contrasts. Across the thesis, productive coverage, lexical diversity, sentence preference, conversational form, and generated grammaticality should consequently be treated as complementary outcomes rather than interchangeable estimates of a single human-likeness score.

\subsection{Uncertainty Includes Sampling and Measurement Choice}

Quantitative uncertainty is broader than an error bar around a mean. It includes variation across corpus samples, model initialisations, generated samples, test items, linguistic phenomena, conversations, and human participants. The relevant unit of analysis differs across studies: generated tokens are nested within model runs, test sentences within paradigms, and utterances within conversations. Confidence intervals and hierarchical analyses should respect these dependencies rather than treating every observation as independent.

A second form of uncertainty arises from analytic choices. In \cref{chp2_chap}, estimated vocabulary trajectories depend on the production threshold, decoding temperature, curve family, frequency-band definition, and any extrapolation beyond observed exposure. Sensitivity analyses are therefore part of the substantive result, not an optional robustness appendix. In \cref{chp3_chap}, item-level bootstrap intervals and mixed-effects modelling support the claim that retrieval benefits are larger for low-frequency items, while uncertainty across phenomena and retrieval configurations limits claims of universality. Similar care is required in \cref{chp4_chap}, where automatic evaluators and dialogue metrics introduce their own measurement error and construct assumptions.

The appropriate response is triangulation. A developmental claim becomes stronger when it survives plausible linking functions and appears across measures whose biases differ. When measures disagree, the disagreement should be preserved and explained rather than averaged away. This principle follows directly from the calibration framework in \cref{chap:litreview}: evidence about a learner is always mediated by a task model. Reporting that mediation explicitly allows the thesis to distinguish robust conclusions from conclusions that remain conditional on a threshold, prompt, decoder, retrieval configuration, or evaluator.

\section{From Static Input to Social Interaction}
\label{sec:gd_interaction}

The studies in this thesis broaden the learning environment under consideration, from distributions in static linguistic input, through access to stored instances, to signals derived from social interaction. This progression should not be interpreted as a single model acquiring successively more child-like components. The studies use different models, tasks, and outcome measures. Rather, they examine complementary ways in which information may become available to a learner. \Cref{chp2_chap} evaluates what predictive learning produces under calibrated exposure and production opportunities; \cref{chp3_chap} asks whether direct access to prior instances can compensate for weaknesses in parametric knowledge; and \cref{chp4_chap} investigates both the simulation of interaction and the use of caregiver-derived signals during learning.

\subsection{Interaction as an Object of Simulation}

The dialogue benchmark in \cref{chp4_chap} first treated language models as \emph{simulators}, asking whether they could reproduce observable properties of child--caregiver conversations. Its multi-level design proved important. Depending on the model and prompting condition, generated language approximated aspects of the CHILDES reference at the word and utterance levels, including lexical density, utterance length, and syntactic complexity. The few-shot prompting condition also brought several measures closer to the human reference, particularly for caregiver-like generation. Because the supplied prompt templates include length constraints as well as examples, these results concern the prompting condition as implemented; they do not isolate examples as the sole cause of changes in utterance length. However, this improvement did not extend uniformly to the dialogue level. Under the study's operationalizations, simulated caregivers remained more semantically aligned and less semantically diverse than human caregivers, and discrepancies often increased when two model instances interacted over multiple turns.

These findings show why simulation fidelity cannot be inferred from locally plausible responses. A model may generate an appropriate-looking reply to a human utterance while producing a different interactional regime when its own outputs shape the subsequent context. In multi-turn simulation, small biases can accumulate as each model conditions its next response on the other's output, producing patterns that are not apparent in single-turn evaluation. This is adaptation through the conversation context, without online parameter learning. Because both interlocutors are simulated, the resulting divergence cannot be assigned to one role alone or directly predict performance with a human interlocutor. Interaction is therefore not merely a longer sequence of independently adequate utterances; it is an evolving relation between interlocutors whose behavior is mutually conditioning.

Crucially, simulation fidelity is not evidence of acquisition. The age-conditioned models were already pretrained and were prompted to enact children or caregivers of particular ages. Apparent age-related differences consequently reveal knowledge that the pretrained model can express under a role instruction, not a developmental trajectory caused by accumulating child-like experience. Likewise, the few-shot contrast demonstrates sensitivity to prompt context and instructions, not lasting learning. The benchmark is informative about the behavioral adequacy of a simulator, but it does not establish how children acquire the behaviors being simulated.

\subsection{From Caregiver Responses to Reward Signals}

The reinforcement-learning study in \cref{chp4_chap} addressed a different question by treating models as controlled learners. Small causal language models were first trained on caregiver utterances and then fine-tuned using rewards derived from naturally occurring child--caregiver pairs. The comparison distinguished communicative feedback, structural alignment, semantic contingency, and affective feedback. These signals produced different outcomes. Structural alignment yielded the clearest gains in the grammaticality of generated utterances; communicative feedback produced smaller and less consistent gains, with clarification requests more informative than acknowledgements. Semantic-contingency and affective rewards did not improve grammaticality under the measures used and sometimes reduced it, while also changing utterance length, lexical entropy, and function-word density.

This finding is stronger than an unadjusted difference in sample averages. In the reported mixed-effects analysis, the reward contrasts control for utterance length and non-word rate and include a generation-run intercept. Structural alignment remains positively associated with ChildCG grammaticality, with estimated changes relative to baseline of $+0.077$ and $+0.060$ at the 1M- and 10M-word scales, respectively (Table~\ref{tab:grammar_reward_effects}). The positive direction is also present under the GEC-based generation measure. Thus the result is not explained away by these two measured covariates, although it remains dependent on automatic assessments of grammaticality. The noisier 0.1M condition further shows that the learner's initial language distribution matters for how effectively the reward can guide production.

This manipulation provides a conditional causal result within the computational system: changing the reward changes the learner's production policy. It does not, however, instantiate complete online interaction. Each frozen reward model maps a generated child utterance to a scalar value learned from correlations between child utterances and subsequent caregiver behavior. During fine-tuning, it neither generates the caregiver's response nor conditions its feedback on an evolving dialogue, shared context, or the learner's history. The learner therefore receives a prediction of the feedback that an utterance tends to elicit, rather than a situated response from an adaptive interlocutor. This abstraction is useful because it isolates candidate signals, but it tests the utility of fixed reward proxies, not the full developmental effect of reciprocal social interaction.

The distinction also prevents the two studies in \cref{chp4_chap} from being treated as a single pipeline. The simulated interlocutors from the benchmark did not supply the rewards used in fine-tuning, and the reward-trained learners were not evaluated on the interaction benchmark. Better simulation would not by itself validate a learning signal, just as a reward-induced production change would not establish that the resulting model sustains more human-like dialogue.

\subsection{Production Gains and Grammatical Generalization}

The divergence between free generation and controlled grammatical evaluation is one of the most informative results of the feedback study. Caregiver-derived rewards, especially structural alignment, could increase the proportion of generated utterances classified as grammatical without consistently improving preferences on BLiMP or Zorro. One possible explanation is a redistribution of probability toward familiar or stable utterance types, without corresponding improvement on the grammatical contrasts tested in the benchmarks. This remains an explanation to test, not a demonstrated account of the structural-alignment effect; its persistence after adjustment for length and non-word rate argues against attributing it simply to shorter outputs. Changes in generation therefore support a claim about productive behavior, but not yet a claim about general grammatical knowledge.

The converse qualification is also necessary. Minimal-pair benchmarks sample particular constructions and may not target the errors that occur frequently in the learner's own productions. Their null results do not prove that no grammar-relevant change occurred. Instead, disagreement between the measures localizes the unresolved question: which regularities changed, and how far do they transfer? As argued in \cref{chap:litreview}, a general acquisition claim requires convergence across measures whose linking assumptions are independently justified.

The grammar-targeted topline provides a further comparison: it improves both generated grammaticality and the reported minimal-pair scores. However, its reward model was itself trained on grammaticality labels from BLiMP and Zorro, and the present description does not establish an independent evaluation split. It therefore shows that explicitly grammar-supervised reward optimization can improve the reported outcomes, but cannot serve as an independently validated upper bound on grammatical generalization or prove that the empirical rewards fail solely because they are indirect. This distinction leaves the caregiver-reward contrast intact while limiting what the topline adds to its interpretation.
% AUTHOR CHECK: document the separation between topline reward training, model selection, and final BLiMP/Zorro evaluation items and templates.

\section{Theoretical Implications for Language Acquisition}
\label{sec:gd_theory}

Taken together, the findings support a plural and functionally differentiated account of the information that may contribute to language learning. Predictive exposure, episodic access, and interaction-derived feedback alter a learner at different stages and in different ways. The thesis does not show that these mechanisms form a single developmental architecture, but it clarifies hypotheses that an integrated account would need to distinguish.

\subsection{Complementary Sources of Support}

Predictive learning provides gradual sensitivity to distributional regularities, but the results of \cref{chp2_chap} suggest that, under the production measure used there, this sensitivity remains strongly shaped by lexical frequency. \Cref{chp3_chap} provides a limited computational demonstration that inference-time access to stored instances can selectively improve performance when baseline decisions are less accurate, narrowing the frequency gap on syntactic contrasts without eliminating it. The result is compatible with the general idea that consolidated statistical knowledge and instance-based access can make complementary contributions \citep{mcclelland1995there,kumaran2016learning}. Because the vocabulary and retrieval studies use different models and outcomes, however, retrieval cannot be claimed to explain or repair the expressive-vocabulary trajectory observed in \cref{chp2_chap}.

Feedback introduces a third form of support: it changes which productions are reinforced and thereby updates the learner's parameters, whereas the retrieval intervention changes access to stored examples at inference time. The structural findings in \cref{chp3_chap,chp4_chap} offer a cautious point of contact. Retrieved instances containing useful structural information supported performance on syntactic contrasts, while a reward associated with caregiver reuse of POS-bigram structure improved generated grammaticality. These manipulations are not equivalent, and neither identifies the mechanism used by children. Together they identify two distinct routes through which structure-related information can improve model behavior: selecting useful instances at inference and reinforcing a production policy during fine-tuning. This motivates a testable hypothesis that the way a learner accesses and uses structural evidence matters alongside the amount of exposure. Whether children exploit analogous information remains an empirical question.

\subsection{Feedback Is Not a Unitary Developmental Resource}

The reward results argue against treating ``caregiver feedback'' as a homogeneous source of evidence. A clarification request, structural reuse, semantic continuation, and positive affect may all be contingent on a child's contribution, yet they need not provide information about the same target. Structural alignment may preserve a reusable linguistic frame; clarification requests may indicate communicative difficulty without locating its source; semantic contingency may support common ground; and affect may regulate engagement rather than grammatical form. Their value should consequently be assessed relative to a specified developmental outcome.

For the same reason, the weaker grammaticality observed under semantic and affective rewards does not imply that such responses are harmful to children. Grammar was the principal outcome of the study, whereas these signals may serve semantic, pragmatic, motivational, or relational functions that were not measured. The theoretically relevant question is therefore not whether feedback is globally beneficial, but which signal can support which change, for which learner, under which interactional conditions.

\subsection{From Computational Effects to Developmental Claims}

Across these studies, computational interventions establish forms of conditional sufficiency rather than mechanistic identity. Retrieval shows that external instance access \emph{can} reduce a model's frequency sensitivity; reward fine-tuning shows that caregiver-correlated scalar signals \emph{can} reshape production. Neither result shows that children use the same representation, optimization procedure, or learning dynamics. In particular, natural interaction supplies temporally situated linguistic responses, multimodal context, joint attention, and mutual adaptation that are absent from a frozen reward function.

The broader theoretical contribution is therefore a disciplined decomposition of the acquisition problem. Distributional exposure determines what regularities are available for gradual learning; memory determines which particular experiences can be re-accessed; and interaction can make some learner behaviors more consequential than others. Their effects must be tested separately before they are combined. Consistent with the calibration framework in \cref{chap:litreview}, claims about development should remain tied to the manipulated information, the learner in which it was manipulated, and the outcome through which its effect was observed.

\section{Thesis-wide Limitations}
\label{sec:gd_limitations}

The main limits concern the available experience, the validity of the measurements, and the populations to which the results can generalize. The distinctions between model roles and developmental scales discussed above also prevent treating the four studies as a cumulative experiment on a single learner.

\subsection{Textual Input and Missing Grounding}

All studies represent language as text. Children encounter speech embedded in perceptual and social events, with acoustic variability, prosody, gaze, gesture, and visible referents \citep{Tomasello1995,Csibra2009}. Transcription removes much of this information while supplying discrete orthographic units. In \cref{chp2_chap}, character input retains spaces and punctuation, so it does not recreate segmentation from continuous speech. In \cref{chp3_chap}, stored text preserves linguistic instances without the perceptual or social content of an experienced episode. In \cref{chp4_chap}, dialogue is represented by written turns, and feedback learning further reduces the caregiver response to a scalar prediction. These simplifications support tractable interventions, but they do not establish what adding any particular missing modality would change. That contribution requires a matched intervention.

\subsection{Automatic Measures and Validation Data}

The thesis makes its task models explicit, but several remain incompletely validated. The \meCDI{} threshold is chosen through correspondence with the human vocabulary trajectory, making independent calibration important for subsequent comparisons. Dictionary-based non-word judgments can penalize legitimate child forms and do not isolate articulation. The dialogue measures quantify linguistic form and embedding-based relations, without fully assessing whether a caregiver addresses a child's intended meaning or repairs a misunderstanding. Their exact definitions, normalization, and uncertainty matter to the interpretation of alignment and diversity.

In the feedback study, automatic annotation, reward prediction, and generated-utterance evaluation introduce separate sources of error. A predictor fitted to corpus utterances may become unreliable on the learner's evolving output distribution. The agreement of ChildCG and GEC on the principal reward pattern is useful converging evidence, but neither substitutes for a targeted human assessment of the affected errors. Likewise, a held-out reward-validation set does not by itself establish separation from the final acquisition test; the grammar topline requires the specific data-provenance check described above.

\subsection{Information Budgets and Population Generalizability}

Unique corpus size, repeated exposure, inference-time memory, and pretrained reward components contribute different information to the experimental system. Their provenance is necessary for claims about learning from child-comparable experience alone. In particular, the smaller retrieval condition combines several text genres and repeated epochs; broad pretraining supports the dialogue simulators and parts of the feedback pipeline. These choices are compatible with controlled computational comparisons, but should not be compressed into a single label of developmental realism.

The studies also concentrate on English and rely heavily on CHILDES as input or reference \citep{macwhinney2000childes}. Recorded families, activities, age ranges, and transcription conventions sample only part of children's experience. Aggregate curves can conceal variation across children and communities \citep{WeislederFernald2013,bergelson2023everyday}. Conclusions about lexical frequency, word classes, and caregiver response may depend on English morphology and word order as well as the sampled interactional practices. Their scope is the populations and corpora analyzed here until those dependencies are tested directly.

\section{Toward Integrated Developmental Models}
\label{sec:gd_future}

The next step is to test the unresolved links between the present results. Integration will be informative when it connects outcomes and resources while keeping their individual contributions identifiable.

\subsection{Longitudinal and Multimodal Experience}

To distinguish differences in vocabulary-growth dynamics from artifacts of independent exposure conditions, a learner could update continuously from a temporally ordered stream. The same learner could be evaluated at predefined checkpoints while its unique exposure, repeated exposure, optimization history, and memory state are recorded separately. Matched text, speech, and grounded conditions could then test whether particular perceptual or referential information changes rare-word availability and the content/function-word balance at the same exposure budget. Such a model would still require an explicit account of what was available to the learner's attention; adding synchronized modalities should not be equated automatically with reproducing a child's effective experience.

\subsection{Adaptive Episodic Memory}

The residual frequency gap and the structural-match analyses motivate memories that are updated online and selected according to more than embedding similarity. Candidate mechanisms include preferential storage of surprising or informative events, retrieval conditioned on structural demands, adaptive context and neighbour selection, interference, forgetting, and gradual consolidation into parametric knowledge. These mechanisms should be compared through controlled ablations, keeping the stored material and target tasks comparable. Applying the intervention to both \meCDI{} generation and syntactic preferences would directly test whether the observed compensation extends across the two frequency-related findings. Their value would lie in explaining when access to prior instances improves transfer, not in treating a successful retrieval system as direct evidence for a particular neural implementation.

\subsection{A Contingent Caregiver}

An interactive learner could produce an utterance, receive a context-sensitive caregiver response, and update before the next exchange. The caregiver would need to condition on dialogue history and the learner's changing behavior, rather than returning a fixed scalar prediction. Its responses should first be validated against natural interaction at utterance and dialogue levels. Yoked, delayed, or shuffled-feedback controls could then isolate contingency while preserving the amount and surface form of input. A factorial comparison of structural and semantic rewards would also test whether their isolated effects combine, trade off, or depend on the learner's current production distribution; the present single-reward comparisons cannot decide this. Because two learned agents can amplify one another's errors, simulation fidelity and learning effects would still require separate evaluation.

\subsection{Held-out, Multi-outcome Evaluation}

Evaluation should follow the learner longitudinally across expressive vocabulary, word--referent mapping, controlled syntactic contrasts, spontaneous production, and sustained interaction. Calibration materials should be separated from final tests, and reward models or evaluators should not be trained on the same items used to support a generalization claim. Targeted contrasts should first be built around recurring errors in both child and model production, such as omissions, and then tested with held-out lexical items and combinations. Additional transfer tests can hold out constructions, speakers, and interaction contexts. Reporting endpoints, curve shapes, item-level errors, and uncertainty together would show whether an intervention produces a narrow production shift or a regularity that generalizes across outcomes.

\subsection{Cross-linguistic and Individual-level Validation}

An integrated account should also be tested across languages and across individual learning histories. The goal is not to make every model reproduce an average child curve, but to ask whether variation in a learner's own experience predicts variation in its subsequent behavior under comparable measures. Multilingual and community-diverse corpora would test whether conclusions about frequency, memory, and feedback survive changes in morphology, word order, interactional practice, and exposure distribution. Wherever privacy and sample size permit, hierarchical analyses should preserve child-, caregiver-, conversation-, item-, and model-level variation rather than collapsing them into one population mean.

\section{General Conclusion}
\label{sec:gd_conclusion}

This thesis asked how limited and uneven linguistic experience is transformed into expressive and grammatical behavior, and how models can clarify the contributions of distributional exposure, access to prior instances, and caregiver-derived feedback. The findings show that the amount of exposure is only part of the explanation: the accessibility of particular instances, the content of a learning signal, and the task used to observe behavior also matter. Developmental interpretation requires the learner's experience, the learner and intervention, the observed outcome, and the linking assumptions to be specified together. Similar performance under one task is not sufficient to establish a shared mechanism, while failure under a mismatched comparison is not sufficient to establish that an outcome is unlearnable.

Within these boundaries, the empirical chapters identify a coherent pattern. The expressive-vocabulary study reveals developmental and frequency-related discrepancies under a production-based comparison. The retrieval study shows that access to stored instances can selectively improve performance on low-frequency syntactic items, although it does not eliminate the frequency gap. The interaction studies show that plausible individual utterances do not guarantee realistic sustained dialogue, and that caregiver-derived rewards have differentiated effects: structural alignment most consistently improves generated grammaticality, without establishing broader grammatical transfer.

The resulting contribution is a set of tools and findings about how linguistic evidence becomes available in behavior: a calibrated expressive-vocabulary measure, an analysis of frequency-selective retrieval, a benchmark of sustained dialogue, and a comparison of caregiver-derived learning signals. Predictive learning, instance access, and feedback make distinct, task-dependent contributions whose effects must be measured separately. Language models are most informative here not as artificial children, but as controlled systems in which experience, mechanisms, and evaluation can be varied. Used with calibrated comparisons and explicit inferential limits, they can help transform broad claims about language acquisition into testable questions about which learning resources are sufficient for which developmental outcomes.
